\section{Introduction}\label{sec:intro} 




% \yd{expensive - mitigate expensive existing works are not good - our observations and this observation gives us a very unique and interesting opportunity (i.e., to local estimate it) but directly use is not good - we propose a framework which is wonderful}



Large language models (LLMs) such as the GPT series~\citep{radford2019language, brown2020language, achiam2023gpt}, LLaMA~\citep{touvron2023llama}, Claude~\citep{anthropic_introducing_claude_2023}, and DeepSeek~\citep{guo2025deepseek} have demonstrated strong performance in real-world machine-learning-as-a-service (MLaaS) applications, ranging from chat assistants to code generation~\citep{park2023thinking, dong2023towards, liu2024exploring, li2026llmclinicalgraphstructure, yu2026health}. As LLMs are increasingly deployed interactively, \textit{multi-turn LLM serving}, which covers extended interactions between humans and LLMs or among LLM agents, has become a key research focus~\citep{yi2024survey, li2025beyond, zhao2025amplifying}.
However, supporting efficient context-dependent 
% \yd{is this word some terminology in this domain? why do we need to put it here?} 
% \xq{yes, they mean we need whole history to process the later turns for LLM, but since I mention stateless later, I will just remove it}
multi-turn interaction remains a key challenge, as most LLM serving systems are stateless and require resending the entire conversation history, including all prior queries and responses, with each new query to generate a new response~\citep{ananda2025beyond, moon2025accelerating}. This leads to redundant computation, high latency, and rising serving costs as conversations lengthen.

Previous work has explored efficient multi-turn LLM serving through two main approaches. One line of work focuses on \textit{single-model methods} that compress dialogue history~\citep{wang2025recursively, chen2024compress, xiao2024efficient}, retrieve memory from external modules~\citep{melz2023enhancing, gutierrez2024hipporag}, or reuse attention computations~\citep{gao2024cost, jeong2025accelerating, anthropic_prompt_caching_2024}. However, these still rely heavily on large LLMs for every turn, leading to high monetary cost, latency, and GPU usage. They also often truncate or overlook extended context, limiting reasoning over long dialogues. Another line adopts \textit{multi-model methods}, routing simple queries to smaller models while escalating harder ones to larger LLMs~\citep{behera2025towards, schick2023toolformer, ding2024hybrid, shnitzer2023large}. Yet, small models struggle to generalize across dialogue complexity, and model switching introduces additional overhead. Moreover, LLMs are known to rely on early turns~\citep{xiao2024efficient, laban2025llms}, compounding the difficulty of maintaining coherence. 
% \yd{I think we did involve both types of baselines into our experiments, right? make sure the answer is yes and do a call back in experimental setup when introducing baselines like ``there are mainly two mainstreams of models xxx" (of course also mention your straightforward methods there as well)} Recognizing these limitations, we ask: 
%\xq{sure}
% \yd{i changed the ``design" to ``achieve" since only designing something is still NOT what people want.}\xq{sure}
% \yd{your question should be whether something is possible, not whether ``you/we" can achieve xx. People care about what the nature allows, not what ``you/we" can do.}

\textit{\textbf{Is it possible to achieve an efficient, context-aware multi-turn LLM serving framework that avoids recomputing the full history at every turn while preserving response quality?}}


% \yd{mention to achieve this goal, we perform in-depth explorations, and it reveals xxx. Plus, I think the logic in this paragraph is unnecessarily complex. follow: to handle gap, we explored more - reveal phenomenon - this phenomenon really makes sense with real-world examples - intuitively, directly dealing with the short turns with a small model can improve efficiency - but we still need to make it context-relevant, i.e., to make it aware of the context contained in the big head - mention existing strategies and mention they are not good and mention this is really difficult}
To answer this question, we examine real-world multi-turn dialogues across domains and find a long-tail distribution of token counts over turns: early turns are substantially longer, while later turns shrink. This phenomenon aligns with the intuition of prior work that early turns typically carry \rev{the substantive openings that set the issues and anchors of a conversation}, including questions, requests, and proposals, while later turns are typically minimal acknowledgments~\citep{he2018decoupling, stolcke2000dialogue}.
% \yd{mention this phenomenon aligns with our intuition of how people typically interacting with the model} 
This long-tail trend gives rise to an intuitive idea: since later turns are relatively short, a smaller language model might suffice to generate responses more cost-efficiently. 
% \yd{starting from here, we involved some irrelevant redundant logic. simply say that the problem comes from the big head: if small models still need to process the head, efficiency is still a problem.} 
The bottleneck, however, is the “big head”: most grounding is established early, and later turns remain highly dependent on it. A small model may behave reasonably at the start, but its responses drift from the large model as the dialogue progresses, so simply handing later turns to a small model degrades quality. The small model must therefore not only process shorter inputs but also approximate the larger model's behavior within the local region of its output space that the prior dialogue induces. We call this a local manifold approximation problem.

Building on these insights, we present SOMA (\textit{S}oft-prompts for l\textit{O}cal \textit{M}anifold \textit{A}pproximation), which adapts a small language model to the local behavior of a large one, conditioned on the early turns of a session, through a three-stage pipeline: (1) \textit{soft-prompt tuning}, which explores the local region induced by the early context to find the directions of maximal behavioral divergence between the two models; (2) \textit{localized fine-tuning}, which aligns the small model with the large one on a handful of mined input--output pairs; and (3) \textit{efficient inference}, which serves later turns from a compressed context and rolls back to the large model when the topic drifts. Our contributions are: 
% \yd{your contributions are too detailed. do this: 1 - Explorations revealing the phenomenon of xxx; 2 - Novel framework for xxx; 3 - Theoretical analysis of xxx; 3 - empirical evaluation for xxx. 4 in total.}
% \yd{name each with textbf bold font. make all exactly 2-rows or 3-rows}
\begin{itemize}[left=0pt,itemsep=2pt]
\item \textbf{Long-tail pattern in multi-turn dialogues.} 
We reveal an under-explored long-tail pattern in multi-turn dialogues: the first few turns concentrate heavy context, while later turns are shorter yet more dependent on earlier ones. This suggests large savings if a small, cheap model can serve the later turns once it is given the accumulated context.

% (The separate "Empirical characterization" point is merged above per feedback.)

% \item \textbf{Empirical characterization.} \yd{we shall remove this -- merge this one into the previous one very briefly would be good enough}
% We quantify tail heaviness across multiple datasets and show strong context reliance as conversations progress. The measurements reveal that a small language model could potentially replace a large language model during the conversation to serve later turns. This empirical characterization suggests that a small model can replace a 

\item \textbf{SOMA: Efficient multi-turn serving.} It first learns soft prompts that expose the largest surrogate–original response dissimilarity, then adapts the surrogate localized LoRA accordingly, enabling prompt-free inference with a simple cosine gate for switching and rollback.

\item \textbf{Theory analysis and empirical evaluation.} 
We provide concentration-based bounds for switching, coverage guarantees for prompt-direction search, and suboptimality limits for selected directions. Guided by these results, empirical studies show the effectiveness of SOMA in real world.
\end{itemize}


% \yd{no need to do this any more. We shall remove this paragraph. We are tight on page limit.}
% The paper is structured as follows: Section~\ref{sec:prelim} introduces preliminaries and empirically characterizes the long-tail token–turn pattern in multi-turn dialogues. Section~\ref{sec:method} presents details of SOMA, including soft-prompt mining to expose least-aligned directions, localized LoRA fine-tuning to adapt the surrogate, and efficiency-aware inference with switching and rollback. Section~\ref{sec:theory} develops the theoretical foundations, including warm-start generalization and detection bounds for switching as well as coverage guarantees for prompt-direction search. Section~\ref{sec:exp} reports empirical experiment results across six datasets. Section~\ref{sec:conclusion} concludes the work.






